% 数学ノート集・共通テンプレート v2（積分論2026年度第1部の版面を参考）
% A5判。章内で定義・定理・例・注意等の番号を共通化。
% コンパイル： platex -interaction=nonstopmode math_note_template_v2.tex (2回)
%              dvipdfmx math_note_template_v2.dvi
\documentclass[a5paper,10pt,onecolumn]{jsarticle}
\usepackage[dvipdfmx,top=22mm,bottom=20mm,inner=19mm,outer=18mm,headheight=13pt,headsep=6mm]{geometry}
\usepackage{amsmath,amssymb,amsthm,mathtools}
\usepackage{enumitem,fancyhdr,titlesec}
\usepackage[dvipdfmx,bookmarks=true,colorlinks=false]{hyperref}
\usepackage{pxjahyper}

% ----- 差し替える文書情報 -----
\newcommand{\NoteTitle}{Picard--Vessiot拡大と微分Galois群}
\newcommand{\NoteSeries}{数学ノート集}

% ----- 文字・組版 -----
% 参考PDFの質素な明朝＋Computer Modernの印象を引き継ぎながら，
% A5判で窮屈にならないよう本文の行間を少し確保する．
\linespread{1.13}
\setlength{\parindent}{1zw}
\setlength{\parskip}{0pt}
\setlength{\abovedisplayskip}{6pt plus 2pt minus 2pt}
\setlength{\belowdisplayskip}{6pt plus 2pt minus 2pt}
\setlength{\abovedisplayshortskip}{3pt plus 2pt minus 1pt}
\setlength{\belowdisplayshortskip}{4pt plus 2pt minus 1pt}
\setlist[enumerate]{topsep=3pt,itemsep=1pt,parsep=0pt,leftmargin=2.4zw}
\titleformat{\section}{\normalfont\normalsize\bfseries}{\thesection}{0.8zw}{}
\titlespacing*{\section}{0pt}{2.0ex plus .4ex}{.9ex}
\titleformat{\subsection}{\normalfont\normalsize\bfseries}{\thesubsection}{0.7zw}{}
\titlespacing*{\subsection}{0pt}{1.5ex plus .3ex}{.6ex}
\raggedbottom

% ----- 柱とページ番号：ページ上部、罫線なし -----
\fancypagestyle{mathnote}{%
  \fancyhf{}%
  \fancyhead[L]{\footnotesize \NoteSeries\ （\NoteTitle）}%
  \fancyhead[R]{\footnotesize \thepage}%
  \renewcommand{\headrulewidth}{0pt}%
  \renewcommand{\footrulewidth}{0pt}%
}
\pagestyle{mathnote}

% ----- 節内通し番号の定義・定理群 -----
\newtheoremstyle{mathnotetheorem}
 {4.5pt}{4.5pt}{\normalfont}{}%
 {\bfseries}{}{0.65zw}%
 {\thmname{#1}\thmnumber{ #2}\thmnote{（#3）}}
\theoremstyle{mathnotetheorem}
\newtheorem{theorem}{定理}[section]
\newtheorem{definition}[theorem]{定義}
\newtheorem{lemma}[theorem]{補題}
\newtheorem{proposition}[theorem]{命題}
\newtheorem{corollary}[theorem]{系}
\newtheorem{example}[theorem]{例}
\newtheorem{remark}[theorem]{注意}
\renewcommand{\proofname}{証明}
\renewcommand{\qedsymbol}{$\square$}
\makeatletter
\renewenvironment{proof}[1][\proofname]{%
  \par\pushQED{\qed}\normalfont\topsep4\p@\@plus2\p@\relax
  \trivlist\item[\hskip\labelsep\bfseries #1]\ignorespaces
}{\popQED\endtrivlist\@endpefalse}
\makeatother

% ----- 使うノートで適宜拡張する記号 -----
\newcommand{\R}{\mathbb{R}}
\newcommand{\norm}[1]{\lVert #1\rVert}
\newcommand{\abs}[1]{\lvert #1\rvert}

\begin{document}
\section*{読了後にできること}
本稿を読むための前提は，微分体，Picard--Vessiot拡大，定数体．
本稿の具体的な到達課題は，一次線形微分方程式のPicard--Vessiot拡大とGalois群を計算する．


\par\noindent\textbf{記号と約束.}\quad
$\operatorname{Spec}A$ は可換環 $A$ の素イデアル全体の集合とする．
\par

\section{微分体と線形方程式}

本稿では標数0の微分体 $(K,\partial)$ を考える．
定数体
\[
 k=K^\partial=\{c\in K:\partial c=0\}
\]
は代数閉体と仮定する．
多項式環，テンソル積，代数閉体上の線形代数，
有限型代数の基本事項を前提とする．

\begin{definition}
環 $R$ の導分とは
\[
 \partial(a+b)=\partial a+\partial b,\qquad
 \partial(ab)=(\partial a)b+a(\partial b)
\]
を満たす写像をいう．
導分を持つ環を微分環という．
$\partial(I)\subset I$ となるイデアルを微分イデアルと呼ぶ．
零と単位イデアル以外に微分イデアルを持たない
微分環を微分単純という．
\end{definition}

\begin{definition}
$A\in M_n(K)$ とし，
\[
 \partial Y=AY
 \tag{1}
\]
を考える．
微分体拡大 $L/K$ における行列
$Y\in\operatorname{GL}_n(L)$ が式 (1) を満たすとき，
$Y$ を基本解行列という．
\end{definition}

\begin{lemma}[基本解行列の比]\label{lem:constant-change}
$Y,Z\in\operatorname{GL}_n(L)$ が共に式 (1) を満たすなら
\[
 Z=YC
\]
となる $C\in\operatorname{GL}_n(L^\partial)$ が存在する．
\end{lemma}
\begin{proof}
$C=Y^{-1}Z$ と置く．
$\partial(Y^{-1})=-Y^{-1}(\partial Y)Y^{-1}$
を用いると
\[
 \partial C=-Y^{-1}AYY^{-1}Z+Y^{-1}AZ=0.
\]
従って各成分は定数であり，$C$ は可逆なので
$C\in\operatorname{GL}_n(L^\partial)$ である．
\end{proof}

\section{Picard--Vessiot環と拡大}

\begin{definition}
式 (1) に対するPicard--Vessiot環とは，
次を満たす微分 $K$-代数 $R$ をいう．
\begin{enumerate}
\item $R$ は微分単純である．
\item $Y\in\operatorname{GL}_n(R)$ が存在して
 $\partial Y=AY$ を満たす．
\item $R=K[Y_{ij},\det(Y)^{-1}]$ である．
\item $R^\partial=k$ である．
\end{enumerate}
$R$ の分数体 $L=\operatorname{Frac}R$ を
Picard--Vessiot拡大という．
\end{definition}

\begin{lemma}[微分イデアルの根基]\label{lem:radical}
標数0の微分環 $S$ で微分イデアル $I$ が与えられたとき，
$\sqrt I$ も微分イデアルである．
\end{lemma}
\begin{proof}
$S/I$ に移して $I=0$ の場合を示せばよい．
$a$ が冪零元で $a^m=0$ とする．
$\partial^m(a^m)$ をLeibniz則で展開すると，
全ての因子にちょうど一回ずつ微分を作用させた項は
$m!(\partial a)^m$ であり，
他の項は少なくとも一つ微分されない $a$ を含む．
したがって
\[
 m!(\partial a)^m\in aS.
\]
これだけでは冪零性は直ちには従わないので，
元の等式 $a^m=0$ から得られる
$ma^{m-1}\partial a=0$ を繰り返し微分し，
帰納的に
$a^{m-j}(\partial a)^{2j-1}=0$
（$1\le j\le m$）を得る．
特に $(\partial a)^{2m-1}=0$ であり，
$\partial a$ は冪零である．
よって零イデアルの根基は微分で保たれる．
一般の $I$ についても商に移せば従う．
\end{proof}

\begin{lemma}[微分単純環の整域性]\label{lem:domain}
標数0の微分単純な非零可換環 $R$ は整域である．
\end{lemma}
\begin{proof}
補題\ref{lem:radical}により零イデアルの根基は
微分イデアルであり，微分単純性から零である．
従って $R$ は被約である．
極小素イデアル $\mathfrak p$ を取る．
$\mathfrak p$ で局所化すると，
被約性と極小性によって
$\mathfrak p R_{\mathfrak p}$ は局所化環の零根基であり，
従って零である．
導分は局所化へ延長されるので，
$a\in\mathfrak p$ に対して
$a/1=0$ なら
$\partial(a)/1=0$ となる．
従って $\partial(a)\in\mathfrak p$ であり，
$\mathfrak p$ は微分イデアルである．
微分単純性により $\mathfrak p=0$ となるので，
零イデアルは素イデアルであり $R$ は整域である．
\end{proof}

\section{定数体を増やさない構成}

\begin{lemma}[微分代数的定数の代数性]\label{lem:alg-constants}
$R$ を有限生成な微分単純 $K$-代数とする．
$R^\partial$ は体である．
また $R^\partial$ の任意の元は $k$ 上代数的である．
\end{lemma}
\begin{proof}
$c\in R^\partial$ が非零なら
$(c)$ は非零微分イデアルである．
微分単純性から $(c)=R$ なので $c$ は単元であり，
その逆元も導分が零である．
従って $R^\partial$ は体となる．

$c\in R^\partial$ が $K$ 上代数的なら，
$K[X]$ における $c$ のモニック最小多項式 $P$ を取る．
$P(c)=0$ を微分すると
$(\partial P)(c)=0$ であり，
$\partial P$ の次数は $\deg P$ より小さい．
従って $\partial P=0$ である．
全係数が $k$ に属するので，$c$ は $k$ 上代数的である．

逆に $c$ が $K$ 上超越的と仮定する．
$R$ は有限生成 $K$-整域であるから，
\[
 R=K[c,z_1,\dots,z_m]/I
\]
という有限型代数の表示を取れる．
$c$ が超越的なので $I\cap K[c]=0$ であり，
$\operatorname{Spec}R\to\operatorname{Spec}K[c]$
は優勢な有限型射である．
有限型代数の特殊化補題によって，
有限個を除く $a\in k$ で
$R/(c-a)\ne0$ となる．
この特殊化補題は，
$K(c)$ 上の有限生成イデアルのGr\"obner基底を取り，
そこに現れる有限個の分母・先頭係数を避けて
$c=a$ と代入すると，商環が零にならないことから得られる．
$k$ は標数0なので無限体であり，
そのような $a\in k$ を選べる．
一方 $c-a$ は非零の定数なので単元であり，
$R/(c-a)=0$ となって矛盾する．
従って $c$ は $K$ 上代数的，従って $k$ 上代数的である．
\end{proof}

\begin{corollary}[定数体の保存]\label{cor:constants}
$k$ が代数閉体なら，
補題\ref{lem:alg-constants}の $R$ について
$R^\partial=k$ となる．
さらに $L=\operatorname{Frac}R$ の定数体も
$L^\partial=k$ である．
\end{corollary}
\begin{proof}
補題\ref{lem:alg-constants}から
$R^\partial$ の元は $k$ 上代数的なので，
$k$ の代数閉性から $R^\partial=k$ となる．
$b\in L^\partial$ を取る．
\[
 J=\{r\in R:rb\in R\}
\]
と置くと，分母を一つ取れば $J$ は非零である．
$rb\in R$ なら
$\partial(rb)=(\partial r)b$ であるため
$\partial r\in J$ となり，$J$ は微分イデアルである．
微分単純性から $J=R$ だから $b\in R$ である．
よって $b\in R^\partial=k$ であり，
$L^\partial=k$ となる．
\end{proof}

\begin{theorem}[Picard--Vessiot拡大の存在]
\label{thm:pv-existence}
$k$ を代数閉体とする．
任意の $A\in M_n(K)$ に対し
$\partial Y=AY$ のPicard--Vessiot拡大が存在する．
\end{theorem}
\begin{proof}
不定元行列 $X=(X_{ij})$ を取り，
\[
 S=K[X_{ij},\det(X)^{-1}]
\]
に導分を $\partial X=AX$ によって延長する．
Leibniz則と
\[
 \partial(\det X)=(\operatorname{tr}A)\det X
\]
によって $\det(X)^{-1}$ の導分も一意に定まる．
$S$ の真の微分イデアル全体は零イデアルを含むので空でない．
鎖の和集合は真の微分イデアルであるから，
Zornの補題で極大な真の微分イデアル $I$ を取れる．
$R=S/I$ は微分単純である．
補題\ref{lem:domain}から $R$ は整域である．
$Y=X\bmod I$ は可逆で
$\partial Y=AY$ を満たし，
$R=K[Y,\det(Y)^{-1}]$ である．
$R$ は有限生成 $K$-代数なので
系\ref{cor:constants}から
$R^\partial=(\operatorname{Frac}R)^\partial=k$ である．
従って $L=\operatorname{Frac}R$ は
所望のPicard--Vessiot拡大となる．
\end{proof}

\section{拡大の一意性}

\begin{theorem}[Picard--Vessiot拡大の一意性]
\label{thm:pv-uniqueness}
同じ線形方程式 $\partial Y=AY$ に対する
二つのPicard--Vessiot拡大 $L_1,L_2$ は
$K$ 上の微分体として同型である．
\end{theorem}
\begin{proof}
それぞれのPicard--Vessiot環を $R_1,R_2$ とし，
基本解行列を $Y_1,Y_2$ とする．
微分 $K$-代数
$R_1\otimes_KR_2$ を考える．
これは有限生成 $K$-代数であり，
自然な導分
$\partial(a\otimes b)=\partial a\otimes b+a\otimes\partial b$
を持つ．
極大真微分イデアル $J$ を取り，
$S=(R_1\otimes_KR_2)/J$ と置く．
$S$ は微分単純であり，
補題\ref{lem:domain}から整域，
系\ref{cor:constants}から定数体は $k$ である．
自然な写像 $R_i\to S$ の核は微分イデアルであり，
$1$ を零へ写さないので真である．
$R_i$ の微分単純性から核は零であり，
二つの $R_i$ を $S$ の部分環とみなせる．
$S$ の中で
\[
 C=Y_1^{-1}Y_2
\]
は補題\ref{lem:constant-change}により
$\operatorname{GL}_n(k)$ の行列である．
従って $Y_2=Y_1C$ であり，
$K[Y_1,\det(Y_1)^{-1}]
 =K[Y_2,\det(Y_2)^{-1}]$
が $S$ 内で成立する．
これは $R_1$ と $R_2$ が同じ部分環になることを意味する．
その分数体も同じなので
$L_1\cong_KL_2$ を得る．
\end{proof}

\section{微分Galois群}

\begin{definition}
Picard--Vessiot拡大 $L/K$ に対して
\[
 \operatorname{Gal}^\partial(L/K)
 =\{\sigma\in\operatorname{Aut}(L/K):
 \sigma\partial=\partial\sigma\}
\]
を微分Galois群という．
\end{definition}

\begin{theorem}[行列による忠実な表現]
\label{thm:matrix-group}
$Y$ を基本解行列とする．
各 $\sigma\in\operatorname{Gal}^\partial(L/K)$ に対し
ただ一つの $C_\sigma\in\operatorname{GL}_n(k)$ が存在して
\[
 \sigma(Y)=YC_\sigma.
\]
写像 $\sigma\mapsto C_\sigma$ は
群の単射準同型を与える．
\end{theorem}
\begin{proof}
$\sigma$ は $K$ を固定し導分と可換なので
$\partial(\sigma Y)=\sigma(\partial Y)
 =\sigma(AY)=A\sigma(Y)$ となる．
補題\ref{lem:constant-change}から
$\sigma(Y)=YC_\sigma$ となる
$C_\sigma\in\operatorname{GL}_n(k)$ が存在し，
$Y$ が可逆だから一意である．
$\sigma,\tau$ に対し
\[
 (\sigma\tau)(Y)=\sigma(YC_\tau)
 =\sigma(Y)C_\tau=YC_\sigma C_\tau
\]
なので準同型である．
$C_\sigma=I$ なら $\sigma$ は $Y$ の全成分を固定し，
$K$ も固定するから $R=K[Y,\det(Y)^{-1}]$ を固定する．
従って分数体 $L$ 全体を固定し $\sigma=1$ となる．
よって単射である．
\end{proof}

\section{微分Galois群の代数性}

\begin{lemma}[定数による基底変換]\label{lem:torsor}
$R$ をPicard--Vessiot環とし，
$B=R\otimes_KR$ とする．
$Y_1=Y\otimes1$，$Y_2=1\otimes Y$ と置き，
\[
 C=Y_1^{-1}Y_2\in\operatorname{GL}_n(B)
\]
と定める．
$H=B^\partial$ は $k$ 上有限生成な代数であり，
\[
 H=k[C_{ij},\det(C)^{-1}],
 \qquad R\otimes_kH\cong R\otimes_KR
 \tag{2}
\]
が成立する．
\end{lemma}
\begin{proof}
$\partial C=0$ は補題\ref{lem:constant-change}と同様の計算による．
$B$ は左側の $R$ のコピーと第二の $Y$ で生成されるが，
$Y_2=Y_1C$ なので
$B=R[C_{ij},\det(C)^{-1}]$ である．
$H_0=k[C_{ij},\det(C)^{-1}]$ とすると
積写像
\[
 R\otimes_kH_0\longrightarrow B
\]
は全射である．

単射を示す．
もし $k$ 上線形独立な定数元
$h_1,\dots,h_m\in H_0$ が
$R$ 上線形従属となれば，
$m$ を最小にする関係
$\sum_{i=1}^mr_ih_i=0$ を取れる．
$B=R\otimes_KR$ は $R$-加群として自由なので，
$R$ の非零元を局所化してもこの関係は保持される．
分数体 $\operatorname{Frac}R$ の上で
$r_m\ne0$ として割り，最後の係数を1にする．
微分すると最後の定数項が消え，
$\sum_{i<m}\partial(r_i/r_m)h_i=0$ となる．
最小性から全ての $\partial(r_i/r_m)=0$ である．
系\ref{cor:constants}によって $r_i/r_m\in k$ なので，
元の関係は $k$-線形従属性となり矛盾する．
従って積写像は単射である．

よって $B\cong R\otimes_kH_0$ となる．
$R^\partial=k$ なので，
$k$-ベクトル空間 $H_0$ の基底を一つ固定して
係数ごとに導分を考えると，
$(R\otimes_kH_0)^\partial=H_0$ となる．
従って $H=B^\partial=H_0$ であり (2) を得る．
\end{proof}

\begin{theorem}[線形代数群としてのGalois群]
\label{thm:alg-group}
微分Galois群は $k$ 上のアフィン代数群の
$k$-有理点全体に同型であり，
定理\ref{thm:matrix-group}の表現を通して
$\operatorname{GL}_n(k)$ のZariski閉部分群となる．
\end{theorem}
\begin{proof}
補題\ref{lem:torsor}の $H$ を考える．
三つのコピーの基本解行列 $Y_1,Y_2,Y_3$ について
\[
 Y_1^{-1}Y_3=(Y_1^{-1}Y_2)(Y_2^{-1}Y_3)
\]
が成立する．
この等式は $H$ 上の余積
$\Delta_H:H\to H\otimes_kH$ を与え，
$Y_1=Y_2$ の対角写像は余単位，
二コピーの交換は反転を与える．
行列積の結合性からHopf代数の公理が成立する．
従って $\mathcal G=\operatorname{Spec}H$ は
$k$ 上のアフィン代数群である．

$k$-代数準同型 $\epsilon:H\to k$ を取る．
式 (2) で第二の $R$ から $B$ への埋め込みを考え，
$\operatorname{id}_R\otimes\epsilon$ を適用すると
\[
 Y\longmapsto Y\,(\epsilon(C_{ij}))
\]
という $K$-微分環の自己準同型を得る．
Hopf代数の逆元に対応する $\epsilon$ の逆を用いると，
この写像は自己同型となる．
逆に任意の微分自己同型 $\sigma$ は
$Y\mapsto YC_\sigma$ という形なので，
$H$ の生成元 $C_{ij}$ に $C_\sigma$ の成分を割り当てる
$k$-点を与える．
両構成は互いに逆であり，
$\mathcal G(k)\cong\operatorname{Gal}^\partial(L/K)$ を得る．

$H$ は $\operatorname{GL}_n$ の座標環
$k[T_{ij},\det(T)^{-1}]$ の商なので，
$\mathcal G$ は $\operatorname{GL}_n$ の閉部分群スキームである．
標数0の代数閉体上の有限型群スキームは被約なので，
その $k$-点全体は $\operatorname{GL}_n(k)$ の
Zariski閉部分群をなす．
\end{proof}

\section{微分Galois対応}

\begin{theorem}[微分Galois対応]\label{thm:correspondence}
$L/K$ をPicard--Vessiot拡大，
$G=\operatorname{Gal}^\partial(L/K)$ とする．
中間微分体 $K\subset M\subset L$ と
$G$ のZariski閉部分群 $H$ の間に
\[
 M\longmapsto\operatorname{Gal}^\partial(L/M),
 \qquad H\longmapsto L^H
\]
という互いに逆の包含反転対応がある．
また $M/K$ がPicard--Vessiot拡大であることと
対応する $H$ が $G$ の正規部分群であることは同値であり，
この場合
\[
 \operatorname{Gal}^\partial(M/K)\cong G/H
\]
が成り立つ．
\end{theorem}
\begin{proof}
補題\ref{lem:torsor}の
$R\otimes_KR\cong R\otimes_k k[G]$
はPicard--Vessiot環が $G$-torsorであることを示す．
従って $G$ の閉部分群 $H$ に対して，
右側の作用を $H$ に制限したときの不変元は
$R$ の $H$-不変元に対応し，
分数体に移すと $L^H$ を与える．
逆に中間微分体 $M$ について，
$M$ を固定する微分自己同型の条件は
$L$ の有限個の元に対する代数的等式の共通零点として
記述できるので，その安定化部分群はZariski閉となる．
ここで「有限個」を得るには，
$k[G]$ のNoether性を用いて，
各 $m\in M$ の固定条件から定まるイデアルの和が
有限個の生成元で生成されることを用いる．

対応が互いに逆であることは，
式 (2) の同型を $M$ 上へ基底変換した
忠実平坦降下によって従う．
具体的には
\[
 L\otimes_ML
\]
の二つの埋め込みから来る一致条件
$1\otimes z=z\otimes1$ は
$z\in M$ と同値である．
一方torsor同型によってこの一致条件は，
$M$ を固定する部分群の全ての点に対して
$z$ が不変であることと同値になる．
従って
$L^{\operatorname{Gal}^\partial(L/M)}=M$，
および
$\operatorname{Gal}^\partial(L/L^H)=H$ が成立する．
最後に正規部分群の場合は $G/H$ が
$M$ への作用を通してGalois群を与える．
逆に $M/K$ がPicard--Vessiotなら，
$K$ を固定する微分自己同型は $M$ を保ち，
その核 $\operatorname{Gal}^\partial(L/M)$ は正規である．
\end{proof}

\section{一次方程式の例}

\begin{proposition}[指数型一次方程式]
\label{prop:exp-example}
$a\in K$ とし，$\partial y=ay$ を考える．
基本解 $t\ne0$ のPicard--Vessiot拡大 $L=K(t)$ では
\[
 \sigma(t)=c_\sigma t\qquad(c_\sigma\in k^\times).
\]
$t$ が $K$ 上超越的ならGalois群は
$\mathbb G_m(k)=k^\times$ である．
$t$ が $K$ 上代数的ならGalois群は有限巡回群である．
\end{proposition}
\begin{proof}
基本解の比は定数なので
$\sigma(t)/t\in k^\times$ となる．
$t$ が超越的なら，任意の $c\in k^\times$ に対して
$K(t)\to K(t)$，$t\mapsto ct$ は体自己同型であり，
導分との可換性も
$\partial(ct)=c at=a(ct)$ から従う．
従ってGalois群は $k^\times$ 全体である．
$t$ が代数的なら $L/K$ は有限次なので
微分Galois群は有限である．
$k^\times$ の有限部分群は単位根からなる巡回群なので
結論を得る．
\end{proof}

\begin{proposition}[積分型一次方程式]
\label{prop:int-example}
$a\in K$ とし，$\partial t=a$ となる元を含む
Picard--Vessiot拡大 $L=K(t)$ を考える．
このとき任意の $\sigma$ に対して
$\sigma(t)=t+c_\sigma$（$c_\sigma\in k$）である．
$t$ が超越的ならGalois群は
$\mathbb G_a(k)=(k,+)$ であり，
$t$ が代数的なら $t\in K$ でGalois群は自明である．
\end{proposition}
\begin{proof}
$\partial(\sigma(t)-t)=a-a=0$ なので
$\sigma(t)-t\in k$ である．
$t$ が超越的なら任意の $c\in k$ に対して
$t\mapsto t+c$ は $K$-体自己同型であり，
導分と可換であるからGalois群は $(k,+)$ となる．
$t$ が代数的ならGalois群は有限であるが，
標数0の加法群 $(k,+)$ に非自明な有限部分群はない．
従ってGalois群は自明である．
微分Galois対応から固定体は $K$ なので $t\in K$ となる．
\end{proof}

\begin{example}[二階の定数係数方程式]
$K=\mathbb C(x)$，$\partial=d/dx$ とし
$y''+y=0$ を考える．
基本解は $e^{ix}$ と $e^{-ix}$ であり，
Picard--Vessiot体は $\mathbb C(x,e^{ix})$ である．
その微分Galois群は $\mathbb G_m(\mathbb C)$ である．
\end{example}
\begin{proof}
$z=e^{ix}$ と置くと $z'=iz$，$z^{-1}=e^{-ix}$
なので二つの基本解を含む微分体は
$\mathbb C(x,z)$ である．
$z$ が $\mathbb C(x)$ 上超越的であることを示す．
もし代数関係
$\sum_{j=0}^mr_j(x)e^{ijx}=0$
（$r_m\ne0$，$r_j\in\mathbb C(x)$）があれば，
分母を払って多項式係数の関係が得られる．
$x=-it$，$t\to+\infty$ と置くと
$e^{ijx}=e^{jt}$ なので，
最大の $j=m$ に対応する項を $e^{mt}$ で割れば
$r_m(-it)$ と他の指数減衰項の和が零となる．
非零有理関数 $r_m(-it)$ は高々多項式的に増減し，
指数的に零へ近づくことはできないので矛盾する．
よって $z$ は超越的である．
命題\ref{prop:exp-example}から
Galois群は $\mathbb C^\times$ となる．
\end{proof}

\begin{remark}
定数体が代数閉であるという仮定は
Picard--Vessiot拡大の存在・一意性と
Galois群の $k$-点による記述に本質的である．
代数閉でない定数体では
群スキームやtorsorを区別して扱う必要がある．
\end{remark}

\section{到達課題の完全解答}
\noindent\textbf{指数方程式の微分Galois群．}
\begin{proposition}
$K=\mathbb C(x)$，微分 $d/dx$ とし，
方程式 $y'=y$ を考える．
$L=K(e^x)$ はPicard--Vessiot拡大であり，
\[
\operatorname{Gal}(L/K)\cong\mathbb C^\times
\]
である．
\end{proposition}
\begin{proof}
$e^x$ は非零解であり，
$L$ はこの基本解とその逆元で生成される微分体である．
$e^x$ は $K$ 上超越的である．
実際，代数関係があれば，分母を払って
$\sum_{j=0}^mp_j(x)e^{jx}=0$（$p_m\ne0$）となる．
$x\to+\infty$ で $e^{mx}$ で割ると
$p_m(x)\to0$ を強制し，非零多項式の増大と矛盾する．
$L$ を複素平面上の有理型関数の微分体とみれば，
微分が零の元は局所定数な有理型関数なので
定数体は $\mathbb C$ のままである．
任意の $K$-微分自己同型 $\sigma$ に対し
$\sigma(e^x)'=\sigma(e^x)$ なので，
$\sigma(e^x)/e^x$ の微分は零である．
従って $\sigma(e^x)=ce^x$（$c\in\mathbb C^\times$）．
逆に各 $c\ne0$ について $e^x\mapsto ce^x$ は
$K$ 上の微分体自己同型を定める．
この対応は群同型である．
\end{proof}
\noindent\textbf{到達確認．}
$y'=2y$ の基本解とGalois群を求めよ．
\par\noindent\textbf{解答．}
基本解は $e^{2x}$ であり，
同様の議論から $\mathbb C^\times$ を得る．
独立に生成した微分体は $K(e^{2x})$ である．
\end{document}